\explainedfigure{fig_success_versus_localisation_error_pose_trained_models}{0.97}{Training with position error, and adding a second lane cue}{The proposed model reads the road from a map drawn into its input, so it needs to know where the car is on the map. Here the position estimate is shifted by a constant lateral distance (horizontal axis, metres; a small heading error grows with it). Success over 40 episodes for the proposed model and four models trained with random position errors; bands are 95\% intervals. Dashed lines: the map is the only source of lane position. Solid green and pink lines: a second, independent and noisy lane measurement (lateral offset and heading) is also given to the model. 40 episodes per point.}{Training with position errors alone does not help (dashed lines, 0.25 to 0.28 at 2 metres of error). Adding a second independent lane cue keeps success between 0.875 and 0.95 for all error sizes up to 2 metres (solid lines).}{Robustness to a wrong position estimate comes from \emph{redundancy} (two sources that can disagree), not from noise augmentation. Caution: the second cue is simulated (simulator truth plus noise), so the result shows that the network can exploit redundancy, not that a real lane detector would be accurate enough.}{fig:posetrained}
